To do = remove evaluation framework, from visulation-tool.

I feel like it's redundant and to much. 

We have to update it in the code.

We have to disable the future.

We do want some kind of dataset selection, option, for runs.

---> right now that is the least imporatnt. 

We have to focus on our test set and the final validation.




The primary objective of the visualisation framework is to support validation of the drought impact data generated by the extraction pipeline. The system provides an interactive geospatial interface in which extracted impacts are displayed at the NUTS-3 level. Users can explore spatial and temporal patterns through filters based on impact category, severity, recency, publication date, and model configuration. Selecting a region enables inspection of the associated impact events, including the original evidence span, extracted classification, and source article text. This functionality supports manual verification of model outputs while also enabling exploratory analysis of regional drought impact distributions.

The framework was implemented as an interactive Streamlit application using PyDeck-based choropleth visualisation. Aggregated impact frequencies are mapped onto preprocessed NUTS geometries, where colour intensity represents regional mention density. Additional interface components allow users to inspect individual articles and compare extracted events across models and time periods.

A secondary component of the framework supports human evaluation of extracted labels. In this validation mode, model-generated impact events are compared against manually reviewed entries stored in a solution bank. Users can approve or reject individual extractions, after which the validated outputs are stored for subsequent evaluation and comparison across models. This procedure provides a lightweight mechanism for assessing extraction consistency and identifying systematic classification errors.




% Validation 
% Recommendaiotns





% Building on the geocoded location mentions mapped to NUTS-3 regions, the system provides an interactive geospatial visualization and evaluation environment to analyse drought-impact distributions and assess extraction quality.

% \subsection{Geospatial Aggregation and Mapping}

% Each geocoded mention is assigned to a NUTS-3 region through point-in-polygon matching against simplified administrative boundaries. The final analysis is restricted to NUTS-3 level granularity, as this provides a sufficiently detailed spatial resolution for interpreting localized drought impacts in the Netherlands.

% Region-level impact intensity is computed by aggregating weighted mention counts per region. This prevents multiple mentions within a single article from disproportionately influencing regional scores. The resulting aggregated dataset is joined with the NUTS-3 GeoJSON to construct a choropleth representation of drought-impact density.

% \subsection{Interactive Visualization}

% The visualization layer is implemented using Streamlit and PyDeck. It renders NUTS-3 regions as an interactive choropleth map, where colour intensity represents normalized mention density. Darker shades indicate regions with higher concentrations of reported drought impacts.

% The interface supports filtering by model type, publication time (monthly buckets), severity, and impact classification. This enables both temporal exploration and comparative analysis across different extraction models.

% \subsection{Evaluation Framework}

% To support iterative model improvement, a human-in-the-loop evaluation system is integrated into the dashboard. Extracted drought impacts are compared against a \textit{solution bank} (validated extractions) and a \textit{rejected bank} (incorrect extractions). Each event is uniquely identified using a deterministic hash over key attributes such as location, severity, and recency.

% Unseen model outputs are surfaced for manual validation, where users can accept or reject individual extractions. Accepted events are stored in the solution bank and contribute to downstream evaluation.

% Model performance is quantified by comparing extracted events against these validation sets, yielding counts of true positives, false positives, and pending cases. From this, precision is computed per model, enabling consistent comparison of different LLM configurations on the same dataset.